% =================================================
% Домашняя работа. Урок 3
% =================================================

\homeworktitle{Домашняя работа}{Составление чисел. Задачи на цифры}
\studentline

\begin{routebox}[Как выполнять]
  В каждом задании сначала подчеркните условие о повторении цифр. Перебор
  начинайте со старшего разряда и записывайте варианты в устойчивом порядке.
  Задание 7 отмечено звёздочкой: в нём требуется не выписать все числа, а
  обосновать их количество.
\end{routebox}

\begin{homeworktask}[Цифры могут повторяться]{1}
  Запишите все двузначные числа, в записи которых используются только цифры
  \(2\) и \(6\). Цифры могут повторяться.
  \answerspace{18mm}
\end{homeworktask}

\begin{homeworktask}[Каждая цифра используется один раз]{2}
  Запишите все трёхзначные числа из цифр \(0,4,9\), если цифры не
  повторяются и используются все три цифры.
  \answerspace{22mm}
\end{homeworktask}

\begin{homeworktask}[Наибольшее и наименьшее]{3}
  Из цифр \(0,3,5,8\) составьте без повторений:
  \begin{enumerate}
    \item наибольшее четырёхзначное число;
    \item наименьшее четырёхзначное число.
  \end{enumerate}
  Объясните положение нуля в каждом числе.
  \answerspace{25mm}
\end{homeworktask}

\begin{homeworktask}[Зависимость между цифрами]{4}
  Запишите все двузначные числа, в которых цифра десятков на \(2\) больше
  цифры единиц.
  \begin{hintbox}
    Начните с цифры единиц \(0\) и увеличивайте её на единицу.
  \end{hintbox}
  \answerspace{22mm}
\end{homeworktask}

\begin{homeworktask}[Сумма цифр]{5}
  Сколько существует двузначных чисел, сумма цифр которых равна \(9\)?
  Выпишите эти числа и укажите их количество.
  \answerspace{28mm}
\end{homeworktask}

\begin{homeworktask}[Все перестановки]{6}
  Запишите все трёхзначные числа, составленные из цифр \(1,4,7\) без
  повторений. Организуйте список по цифре сотен.
  \answerspace{28mm}
\end{homeworktask}

\begin{homeworktask}[Задание со звёздочкой]{7*}
  Сколько трёхзначных чисел можно составить из цифр \(0,2,5,7\), если цифры
  не повторяются? Выписывать все числа необязательно, но ход подсчёта нужно
  объяснить.

  \begin{hintbox}
    Для первой цифры ноль запрещён. После выбора первой цифры посчитайте
    варианты второй и третьей цифр.
  \end{hintbox}
  \answerspace{33mm}
\end{homeworktask}

\begin{homeworktask}[Разрядный состав]{8}
  Какое число получится, если взять \(6\) десятков тысяч и \(128\) сотен?
  Запишите вычисление.
  \answerspace{22mm}
\end{homeworktask}

\begin{selfcheck}
  \begin{itemize}[label=\(\square\),leftmargin=1.7em,itemsep=0.25em,topsep=0.2em]
    \item Я учёл, могут ли цифры повторяться.
    \item Ни одно многозначное число не начинается с нуля.
    \item Перебор организован по старшему разряду.
    \item Каждое выписанное число удовлетворяет условию.
    \item Я объяснил, почему других вариантов нет.
  \end{itemize}
\end{selfcheck}
